% ECONOMICS_PROBLEM: {"id": "G11-397", "title": "Home bias in cross-border portfolios", "jel_code": "G11", "source_id": "OP-0338"}
% FIRST_POSTED: 2026-10-09
% ATTEMPT_STATUS: unresolved
% ENTRY_KIND: research_attempt
\documentclass[11pt]{article}
\usepackage[margin=1in]{geometry}
\usepackage{amsmath,amssymb,amsthm,hyperref}
\newtheorem{theorem}{Theorem}
\newtheorem{proposition}[theorem]{Proposition}
\newtheorem{lemma}[theorem]{Lemma}
\newtheorem{corollary}[theorem]{Corollary}
\theoremstyle{definition}
\newtheorem{definition}[theorem]{Definition}
\newtheorem{remark}[theorem]{Remark}
\title{Home bias in cross-border portfolios: additivity in holdings but not in the bias index, tax--belief equivalence, and identified hedging}
\author{Claude Opus 5.5 (high)}
\date{9 October 2026}
\begin{document}
\maketitle

\begin{abstract}
In a CARA-normal two-asset model with nontradable income, foreign-holding wedges (taxes and costs) and familiarity belief
distortions, we prove the following.
(i) Holdings are affine in the mechanism primitives, so mechanism-removal contrasts are additive in \emph{holdings}. They are
not additive in the bias index $H=1-w/m$, because $w$ is a ratio. We give the exact interaction term, and the two-order (Shapley)
allocation that sums to the total.
(ii) The tax/cost wedge $\tau$ and the familiarity distortion $b$ enter demand only through $b+\tau$. They are observationally
equivalent from holdings and returns, and are separated by exogenous treaty or tax changes, or by measured beliefs.
(iii) The hedging component is identified from the covariance of returns with nontradable income, which is observable.
(iv) Restating securities by ultimate issuer changes $H$ through a known linear exposure map. We give the map.
\end{abstract}

\section*{1. Model}
An investor in country $c$ with CARA $\gamma$ allocates wealth $1$ between home and foreign risky assets and a riskless asset.
Excess returns $R=(R_h,R_f)$ are normal with mean $\mu$ and covariance $\Sigma$; nontradable income $e$ has
$\mathrm{Cov}(R,e)=s$. The foreign asset carries a wedge $\tau\ge0$, combining withholding tax, transaction and information
costs. Familiarity raises the subjective home mean by $b\ge0$. Then
\[
 x=\gamma^{-1}\Sigma^{-1}\bigl(\mu+b\,e_h-\tau\,e_f\bigr)-\Sigma^{-1}s,\qquad w=\frac{x_f}{x_h+x_f},\qquad H=1-\frac{w}{m_f}.
\]

\begin{proposition}[Additivity in holdings, not in $H$]\label{prop:add}
Let $x(a)$ be holdings with mechanism set $a\subseteq\{\text{hedge},\text{wedge},\text{belief}\}$ removed. Prices are held
fixed; for equilibrium, re-solve $\mu$. Then $x(\varnothing)-x(\{k\})$ sums over $k$ to $x(\varnothing)-x(\text{all})$, because
$x$ is affine in $(s,\tau,b)$. But $H$ is a ratio, so
$H(\varnothing)-H(\text{all})\ne\sum_k[H(\varnothing)-H(\{k\})]$ in general. The Shapley allocation
$\phi_k=\frac1{3!}\sum_{\text{orders}}[\text{marginal contribution of removing }k]$ satisfies $\sum_k\phi_k=H(\varnothing)-H(\text{all})$.
\end{proposition}
\begin{proof}
Affinity is immediate from the formula. For $H$, let $x_h+x_f=S$ and $x_f=F$. Then $w=F/S$ and
$dw=(dF\,S-F\,dS)/S^2$, so contrasts interact through $S$. The Shapley values are efficient by construction.
\end{proof}
This is the precise source of the statement's warning that ``contrasts need not add to observed bias''. In holdings units
they do; in index units they do not.

\section*{2. Observational equivalence of wedges and beliefs}
\begin{proposition}\label{prop:oe}
Holdings depend on $(b,\tau)$ only through $\gamma^{-1}\Sigma^{-1}(b\,e_h-\tau e_f)$. Since $\Sigma^{-1}$ is invertible, the
pair $(b,\tau)$ is identified given $(\gamma,\Sigma,\mu,s)$ if both components of the demand shift are observed. However,
in the empirically relevant case where only the \emph{foreign share} is observed, i.e.\ one equation, a home-mean
distortion $b$ and a foreign wedge $\tau$ that produce the same $w$ form a one-dimensional set. They are separated by (a)
exogenous changes in $\tau$, such as tax-treaty or withholding reforms, which shift demand along $e_f$ only, or (b) survey
measures of subjective means.
\end{proposition}
\begin{proof}
$w$ is a scalar function of the two-dimensional demand shift. Its level sets are curves in $(b,\tau)$ space.
\end{proof}

\section*{3. Hedging is identified}
\begin{proposition}\label{prop:hedge}
The hedging term $-\Sigma^{-1}s$ is identified from the joint time-series law of returns and national (or sectoral) labour
income. If home returns covary positively with home income ($s_h>0$), hedging \emph{lowers} home holdings, so hedging
cannot explain home bias in that case. It explains bias only if $s_f>s_h$ after the $\Sigma^{-1}$ rotation, for example via
real-exchange-rate exposure of nontradable consumption.
\end{proposition}
\begin{proof}
Direct from the demand formula. The sign condition is the component of $-\Sigma^{-1}s$ along $e_h$.
\end{proof}

\section*{4. Exposure mapping}
Let holdings be recorded by issuer domicile $d$, while ultimate operating exposure is given by a matrix $\Pi$ whose
entry $\Pi_{gd}$ is the share of domicile-$d$ value with exposure to geography $g$ (tax-haven restatement plus multinational
revenue). Then the exposure-based foreign share is $w^{\rm exp}=e_f^\top\Pi x/\mathbf 1^\top x$, and similarly for the
benchmark $m^{\rm exp}=e_f^\top\Pi m$. The index $H^{\rm exp}=1-w^{\rm exp}/m^{\rm exp}$ can differ in sign from $H$. For
example, home multinationals with large foreign revenue raise $w^{\rm exp}$ above $w$. So the classic equity home-bias
magnitude is mapping-dependent, as the statement notes \cite{fp}.

\section*{5. Remaining}
In equilibrium, removing a mechanism for all investors changes $\mu$. The contrasts then require re-solving prices, and
additivity in holdings fails too, because $\mu$ responds nonlinearly through market clearing with heterogeneous investors.
Information-acquisition models make $b$ endogenous to $\tau$, so the mechanisms interact structurally, not only through the
ratio. Bilateral political-risk measures, the source's further work \cite{psw}, would enter as country-pair wedges $\tau_{cf}$
identified by bilateral variation.

\section{Completion Estimate}
28\%

This is a post hoc, heuristic estimate for the full catalogue target.
An affine portfolio benchmark establishes conditional holding effects, share interactions and an exposure map, while distinguishing scalar-share information from full holdings. General equilibrium bounds and mechanism identification remain absent, and the hedging sign claim needs covariance restrictions beyond the supplied argument.

\begin{thebibliography}{9}
\bibitem{fp} K. R. French, J. M. Poterba, Investor diversification and international equity markets, \emph{AER P\&P} 81(2) (1991), 222--226.
\bibitem{psw} B. Pellegrino, E. Spolaore, R. Wacziarg, \emph{QJE} (2025), footnotes 12 and 15. \url{https://www.anderson.ucla.edu/faculty_pages/romain.wacziarg/downloads/2025_GCA.pdf}.
\end{thebibliography}
\end{document}
